\documentclass[../../../include/open-logic-section]{subfiles}

\begin{document}

\olfileid{sth}{z}{story}
\olsection{更详尽的叙述}
	
在\olref[story][approach]{sec}中，我们引用了 Schoenfield 对集合形成过程的描述。现在我们再写下几条原则，以使这一叙述更为精确。如下：

\begin{enumerate}
\item[] \stageshier. 每个集合都在某个阶段形成。
\item[] \stagesord. 阶段是有序的：一些阶段在另一些阶段\emph{之前}。\footnote{我们实际上会默认假设这些阶段是\emph{良序}的。其含义将在\olref[ordinals][]{chap}中解释。这是一个实质性的假设。事实上，利用 \citet{Scott1974} 提出的一项非常巧妙的技巧，可以\emph{避免}这一假设，进而将其\emph{推导}出来。（这也将解释为何我们应当认为存在一个初始阶段。）这里我们无法深入讨论；更多内容请参见 \citet{ButtonLT1}。} 
\item[] \stagesacc. 对任意阶段 $S$，以及在阶段~$S$\emph{之前}形成的任意一些集合：在阶段~$S$ 会形成一个集合，其成员恰好就是那些集合。在阶段~$S$ 不形成其他任何东西。
\end{enumerate}

这些都是非形式的原则，但我们将能够用它们为 Zermelo（策梅洛）集合论的若干公理提供辩护。

（需要提醒一句。尽管我们将给出的公理及其名称都是完全标准的，但我们刚刚提出的这些斜体原则在文献中并没有特定的名称。我们只是给它们起了一些希望能有所帮助的名称。）

\end{document}